% !TITLE 木兰花令·拟古决绝词
% !AUTHOR 纳兰性德
% !DYNASTY 明清
% !GENRE 词
% !ORDER 3

\begin{poem}
人生若只如初见\textsuperscript{1}，何事秋风悲画扇\textsuperscript{2}。
等闲\textsuperscript{3}变却故人心\textsuperscript{4}，却道故人心易变。

骊山语罢清宵半\textsuperscript{5}，泪雨霖铃\textsuperscript{6}终不怨。
何如薄幸\textsuperscript{7}锦衣郎\textsuperscript{8}，比翼连枝\textsuperscript{9}当日愿。
\end{poem}

\begin{annotations}
\notetext{1}{初见：初次相见。如果人与人的关系永远像初见时那样美好——开篇即是一句穿透人心的感慨。}
\notetext{2}{秋风悲画扇：扇子到了秋天就被搁置不用（“秋扇见捐”）。典出汉代班婕妤《怨歌行》，以秋扇比喻被遗弃的女子。}
\notetext{3}{等闲：轻易、随便。轻易就变了心。}
\notetext{4}{故人心：旧人的心意、从前爱人的心。}
\notetext{5}{骊山语罢清宵半：骊山（lí shān），在今陕西临潼。唐玄宗与杨贵妃曾在骊山华清宫长生殿七夕夜半私语，誓愿世世为夫妇。}
\notetext{6}{霖铃：即《雨霖铃》，唐代教坊曲名。相传唐玄宗在安史之乱中逃往蜀地，雨中听到铃声，思念杨贵妃而作此曲。}
\notetext{7}{薄幸：薄情、负心。}
\notetext{8}{锦衣郎：身穿锦衣的男子，指富贵人家的公子，这里借指唐玄宗。}
\notetext{9}{比翼连枝：比翼鸟和连理枝。比翼鸟一日一翼，雌雄并飞；连理枝两棵树的枝条连生在一起。比喻夫妻恩爱永不分离。白居易《长恨歌》：“在天愿作比翼鸟，在地愿为连理枝。”}
\end{annotations}

\begin{appreciation}
这首词第一句"人生若只如初见"，是中国词里流传最广的名句之一。题目里的"决绝"，就是一刀两断的意思；"拟古"，是模仿古人的笔法。纳兰性德替古人写一篇断交词，却写出了千古的心事。

"人生若只如初见"——假如人与人的关系，永远停在第一次见面的那一天。初见的时候什么都是好的，后来的事一件都没有发生。可世上的事偏偏不如此。

"何事秋风悲画扇"——扇子到了秋天就被收起来不用了。这是汉代班婕妤的典故：班婕妤失宠后写诗，把自己比作秋天的扇子。秋风画扇，从此成了"被抛弃"的代名词。这两句连起来：为什么相爱的人总要走到秋扇见捐的那一步？

下面两句最狠："等闲变却故人心，却道故人心易变。"明明是你轻易变了心，反而说是我的心思变了。两句写尽负心人的嘴脸——抢先倒打一耙，好让自己占住理。

后四句用唐玄宗和杨贵妃的事。骊山长生殿，七夕夜半，两人私语，发誓世世为夫妇——"在天愿作比翼鸟，在地愿为连理枝"说的就是这一夜。后来马嵬坡兵变，贵妃赐死；玄宗逃往蜀地，雨中闻铃，作了《雨霖铃》曲。"泪雨霖铃终不怨"——贵妃到死没有一句怨言。可那个锦衣郎呢？一个"薄幸"，把皇帝也骂了：当日的比翼连枝，到头来只是一句空话。

弃妇的题目，《氓》早就写过："信誓旦旦，不思其反。"誓言说得那么诚恳，你转头就不认账。两千五百年后，纳兰再说一遍——世道变了，人心变起来的样子，一点没变。《行行重行行》里的女子还在说"努力加餐饭"，那是还有盼头的等待；到了"人生若只如初见"，已经不等了——这是追悔，不是期待。

我在《氓》那首诗里就叮嘱过你：不是每个人都能遇见关雎、淇奥里那样的君子。这首词是同一句话的另一面：人心是会变的，而且变的人从来不承认是自己变的。你要是将来碰见这样的人，记住这两句——等闲变却故人心，却道故人心易变。别跟他辩，辩不赢的，掉头走人就是了。
\end{appreciation}
